\chapter{Ray and Path Tracing}

During this chapter, we will cover the basics of ray tracing and path tracing, two important techniques in computer graphics for rendering images. In order to introduce them, some previous concepts are required.

\begin{definicion}[Light Ray]
    A \emph{light ray} is a straight line along which light energy (photons) travels. It is defined by its origin and direction (which can also be given as a second point).
\end{definicion}

\begin{definicion}[Light Path]
    A \emph{light path} is a sequence of light rays that represent the journey of light from its source to the observer, including any interactions with objects in the scene (such as reflection, refraction, or absorption).
\end{definicion}

When a light ray interacts with an object, it can be absorved or its direction can be changed due to reflection or refraction. In order to render images, there are two main approaches:
\begin{itemize}
    \item \ul{Going forward}: rays are traced forward from the light source, so the entire scene would be perfectly illuminated. However, most of the rays will not reach the observer, so this method is inefficient.
    \item \ul{Going backward}: rays are traced backward from the observer, so only the rays that reach the observer are considered. A ray is traced for determining the color of each pixel in the image.
\end{itemize}


\section{Ray Tracing}

For each pixel in the image, a \emph{primary ray} is traced from the observer's eye through the pixel into the scene. Once the ray intersects with an object, three type of \emph{secondary rays} can be generated:
\begin{itemize}
    \item \ul{Shadow rays}: For each light source, a shadow ray is traced from the intersection point to the light source to determine if the point is in shadow or not. If the shadow ray intersects another object before reaching the light source, the point is in shadow and does not receive direct illumination from that light source.
    \item \ul{Reflection rays}: If the surface is reflective, a reflection ray is traced in the direction of the reflected ray, which is calculated based on the angle of incidence and the surface normal.
    \item \ul{Refraction rays}: If the surface is translucid, a refraction ray is traced in the direction of the refracted ray, which is calculated based on Snell's law.
\end{itemize}

Therefore, for each intersection $N+2$ rays are traced, where $N$ is the number of light sources in the scene. The reflection and refraction rays are recursively considered as primary rays, so they can generate more secondary rays if they intersect with other objects. This process continues until once of the following conditions is met:





\subsection{Diffuse surfaces}



There are two main types of surfaces:
\begin{itemize}
    \item \ul{Diffuse surfaces}: they reflect light uniformly in all directions.
    \item \ul{Specular surfaces}: they reflect light in a specific direction, following the law of reflection.
\end{itemize}